\documentclass[9pt,aspectratio=169]{beamer}

\usepackage{graphicx}
% \usepackage[OT1, T1]{fontenc}
\usepackage{cmbright}
% \usepackage{sansmath}
\usepackage{amsmath}
\usepackage{mathrsfs}
\usepackage{caption}
\usepackage{bm}
\usepackage{bbm}
%\usepackage{enumitem}
\usepackage{natbib}
\usepackage{booktabs}
\definecolor{theme}{HTML}{2e24b1}
\usefonttheme{professionalfonts}
\DeclareSymbolFont{greekletters}{OML}{sansmath}{m}{it}
\DeclareMathSymbol{\beta}{\mathord}{greekletters}{"0C}
\usepackage{etoolbox}
%\DeclareMathSymbol{\delta}{\mathord}{greekletters}{"0E}
%\setbeamercolor{structure}{fg=theme}
% \usetheme{default}
%\usetheme{AnnArbor}
%\usetheme{Antibes}
%\usetheme{Bergen}
%\usetheme{Berkeley}
%\usetheme{Berlin}
%\usetheme{Boadilla}
%\usetheme{boxes}
%\usetheme{CambridgeUS}
%\usetheme{Copenhagen}
%\usetheme{Darmstadt}
\usetheme{Frankfurt}
%\usetheme{Goettingen}
%\usetheme{Hannover}
%\usetheme{Ilmenau}
%\usetheme{JuanLesPins}
%\usetheme{Luebeck}
%\usetheme{Madrid}
%\usetheme{Malmoe}
%\usetheme{Marburg}
%\usetheme{Montpellier}
%\usetheme{PaloAlto}
%\usetheme{Pittsburgh}
%\usetheme{Rochester}
%\usetheme{Singapore}
%\usetheme{Szeged}
%\usetheme{Warsaw}

% \usecolortheme{default}
% \usecolortheme{albatross}
% \usecolortheme{beaver}
% \usecolortheme{beetle}
% \usecolortheme{crane}
% \usecolortheme{dolphin}
% \usecolortheme{dove}
% \usecolortheme{fly}
% \usecolortheme{lily}
% \usecolortheme{orchid}
% \usecolortheme{rose}
% \usecolortheme{seagull}
\usecolortheme{seahorse}
% \usecolortheme{whale}
% \usecolortheme{wolverine}
\newcommand{\aA}{\mathscr A}
\newcommand{\cC}{\mathcal C}
\newcommand{\dD}{\mathscr D}
\newcommand{\bB}{\mathcal B}
\newcommand{\oO}{\mathcal O}
\newcommand{\gG}{\mathscr G}
\newcommand{\hH}{\mathcal H}
\newcommand{\kK}{\mathcal K}
\newcommand{\iI}{\mathcal I}
\newcommand{\mM}{\mathcal M}
\newcommand{\eE}{\mathcal E}
\newcommand{\fF}{\mathcal F}
\newcommand{\qQ}{\mathcal Q}
\newcommand{\tT}{\mathcal T}
\newcommand{\xX}{\mathcal X}
\newcommand{\yY}{\mathcal Y}
\newcommand{\wW}{\mathcal W}
\newcommand{\uU}{\mathcal U}
\newcommand{\lL}{\mathcal L}
\newcommand{\rR}{\mathcal R}
\newcommand{\zZ}{\mathcal Z}

\newcommand{\sS}{\mathscr S}

\newcommand{\pP}{\mathcal P}

\newcommand{\vV}{\mathcal V}

\newcommand{\BB}{\mathbbm B}
\newcommand{\RR}{\mathbbm R}
\newcommand{\CC}{\mathbbm C}
\newcommand{\QQ}{\mathbbm Q}
\newcommand{\NN}{\mathbbm N}
\newcommand{\GG}{\mathbbm G}
\newcommand{\UU}{\mathbbm U}
\newcommand{\ZZ}{\mathbbm Z}
\newcommand{\KK}{\mathbbm K}
\newcommand{\PP}{\mathbbm P}
\newcommand{\EE}{\mathbbm E}
\newcommand{\TT}{\mathbbm T}

\newcommand{\var}{\mathbbm V}
\renewcommand{\implies}{\Rightarrow}
\renewcommand{\iff}{\Leftrightarrow}

\newcommand{\II}{\mathbb I}
\newcommand{\WW}{\mathbb W}

\renewcommand{\phi}{\varphi}

\newcommand{\XX}{\mathsf X}
\newcommand{\YY}{\mathsf Y}

\newcommand{\bP}{\mathbf P}
\newcommand{\bQ}{\mathbf Q}
\newcommand{\bE}{\mathbf E}
\newcommand{\bM}{\mathbf M}
\newcommand{\bX}{\mathbf X}
\newcommand{\bY}{\mathbf Y}
\newcommand{\interior}{\text{int}\,}

\setbeamertemplate{theorems}[numbered]
% \newtheorem{lemma}[theorem]{Lemma}
\newtheorem{proposition}[theorem]{Proposition}

\renewcommand{\baselinestretch}{1.3}
\setlength{\parskip}{1.0ex plus0.5ex minus0.5ex}
\setbeamertemplate{caption}[numbered]
\setbeamertemplate{frametitle continuation}{}
\setbeamertemplate{headline}{}
\defbeamertemplate{footline}{myframe number}
{
  \hspace{1pt}
    \usebeamercolor[structure]{page number in head/foot}%
  \usebeamerfont{page number in head/foot}%
  \raisebox{1.7ex}[0pt][0pt]{% <--- change here
    \insertframenumber\,/\,\inserttotalframenumber\kern1em}%
}
\setbeamertemplate{footline}[myframe number]

\makeatletter
\patchcmd{\beamer@sectionintoc}
  {\vfill}
  {\vskip\itemsep}
  {}
  {}
\makeatother  

\AtBeginSection[]
{
  \begin{frame}
    \frametitle{Outline}
    \begin{minipage}{1.0\linewidth}
      \tableofcontents[currentsection,currentsubsection]
    \end{minipage}
  \end{frame}
}

\AtBeginSubsection[]
{
  \begin{frame}
    \frametitle{Outline}
    \begin{minipage}{1.0\linewidth}
      \tableofcontents[currentsection,currentsubsection]
    \end{minipage}
  \end{frame}
}


\title[Lecture 6: Analysis]{Lecture 6: Analysis}

\author{Junnan Zhang} % Your name
\institute[Paula and Gregory Chow Institute for Studies in Economics\\
Xiamen University] % Your institution as it will appear on the bottom of every slide, may be shorthand to save space
{
	Paula and Gregory Chow Institute for Studies in Economics \\ Xiamen University  \\ % Your institution for the title page
	\medskip
}
\date{Fall, 2026} % Date, can be changed to a custom date
\setbeamertemplate{itemize items}[default]
\setbeamertemplate{sidebar}[sidebar right]
% \definecolor{myco}{HTML}{a67678}
\definecolor{myco}{HTML}{8776a6}
\setbeamercolor{itemize item}{fg=myco} % all frames will have red bullets
\setbeamercolor{itemize subitem}{fg=myco} % all frames will have red bullets
\setbeamercolor{block title}{bg=myco}
\setbeamercolor{structure}{fg=myco}
% - Either use conference name or its abbreviation.
% - Not really informative to the audience, more for people (including
%   yourself) who are reading the slides online

\subject{Economics}
% This is only inserted into the PDF information catalog. Can be left
% out. 

% If you have a file called "university-logo-filename.xxx", where xxx
% is a graphic format that can be processed by latex or pdflatex,
% resp., then you can add a logo as follows:

% \pgfdeclareimage[height=0.5cm]{university-logo}{university-logo-filename}
% \logo{\pgfuseimage{university-logo}}

% Delete this, if you do not want the table of contents to pop up at
% the beginning of each subsection:
% \AtBeginSubsection[]
% {
%   \begin{frame}<beamer>{Outline}
%     \tableofcontents[currentsection,currentsubsection]
%   \end{frame}
% }

% Let's get started
\begin{document}

\begin{frame}
    \titlepage
\end{frame}

\begin{frame}
    \frametitle{Introduction}
    \begin{itemize}
        \item A central concept in analysis is ``convergence''
        \item Building on convergence, we have continuity, completeness,
        compactness, etc
        \item Important theorems then follow
    \end{itemize}
\end{frame}

\begin{frame}{Summary}
    Topics to be covered today
    \begin{itemize}
        \item Metric Spaces
        \item Sequences and Convergence
        \item Completeness
        \item Compactness
        \item Continuity
    \end{itemize}
\end{frame}

\section{Metric Spaces}

\begin{frame}{Distance}
    \begin{itemize}
        \item What is the distance between two real numbers $a$ and $b$?
        \item What is the distance between two points on the plane? 
        \item What is the distance between $f(x) = 0$ and $f(x) = \sin(x)$?
    \end{itemize}
\end{frame}

\begin{frame}
    \frametitle{Metric Space}
    \begin{definition}
        A \textbf{distance} or \textbf{metric} on a set $X$ is a function $d: X
        \times X \to \RR$ such that for all $x, y, z \in X$ the following
        holds:
        \begin{enumerate}
            \item $d(x, y) \geq 0$ and $d(x, y) = 0$ iff $x = y$
            \item $d(x, y) = d(y, x)$
            \item triangle inequality: $d(x, y) \leq d(x, z) + d(z, y)$
        \end{enumerate}
        A \textbf{metric space} $(X, d)$ is a set $X$ together with a distance
        function $d$ on $X$.
    \end{definition}
\end{frame}

\begin{frame}
    \frametitle{Metric Space}
    \framesubtitle{Examples}
    \begin{itemize}
        \item On $\RR$, $d(x, y) = |x - y|$ is a metric (prove it)
        \item More generally, the \textbf{Euclidean} metric on $\RR^n$:
        \begin{eqnarray*}
            d(\bm{x}, \bm{y}) = |\bm{x} - \bm{y}| = \sqrt{(x_1 - y_1)^2 +
                \ldots + (x_n - y_n)^2}
        \end{eqnarray*}
        \item The metric induced by the \textbf{sup norm} on the set of bounded
        real-valued functions on $X$:
        \begin{equation*}
            d(f, g) = \|f - g\|_\infty = \sup_{x \in X}|f(x) - g(x)|
        \end{equation*}
    \end{itemize}
\end{frame}

\begin{frame}
    \frametitle{Open Balls and Neighborhoods}
    \begin{definition}
        \begin{itemize}
            \item An \textbf{open ball} with center $a$ and radius $r$ is
            defined by
            \begin{equation*}
                B_r(a) := \{x\in X: d(x, a) < r \}
            \end{equation*}
            \item A subset $Y$ of $X$ is a \textbf{neighborhood} of $a \in X$
            if there exists $r > 0$ such that $B_r(a) \subset Y$
            \item A subset $S$ of $X$ is \textbf{bounded} if $S \subset
            B_r(a)$ for some $a\in X$ and $r > 0$
        \end{itemize}
    \end{definition}
\end{frame}

\begin{frame}
    \frametitle{Open Balls and Neighborhoods}
    \framesubtitle{Examples}
    \begin{itemize}
        \item $(0, 2)$ is an open ball $B_1(1)$
        \item Any open interval $(a, b)$ in $\RR$ is an open ball (prove it)
    \end{itemize}
    \begin{proposition}
        If $S$ is a bounded subset of a metric space $X$, then for every $y
        \in X$ there exists $\rho > 0$ such that $S \subset B_\rho(y)$.
    \end{proposition}
    Prove it in the assignment
\end{frame}

\begin{frame}
    \frametitle{Interior, Exterior, Boundary, and Closure}
    \begin{definition}
        Let $S$ be a subset of a metric space $(X, d)$.
        \begin{itemize}
            \item A point $x \in X$ is an \textbf{interior} (resp.
            \textbf{exterior}) point of $S$ if some open ball centered at $x$
            is a subset of $S$ (resp. $S^c$).
            \item If every open ball centered at $x$ contains at least one
            point in $S$ and at least one point in $S^c$, then $x$ is a
            \textbf{boundary} point of $S$.
            \item The set of all interior (exterior) points of $S$ is called
            the interior (exterior) of $S$ and is denoted by $\interior S$
            ($\text{ext}\,S$).
            \item The set of all boundary points of $S$ is called the
            boundary of $S$ and is denoted by $\partial S$.
        \end{itemize}
    \end{definition}
\end{frame}

\begin{frame}
    \frametitle{Interior, Exterior, Boundary, and Closure}
    \framesubtitle{Examples}
    Let $X = \RR$ and $S = [0, 1)$
    \begin{itemize}
        \item $x = 1$, $x = 0$, $x = 0.5$, $x = 2$?
        \item What is $\interior S$? What is $\partial S$?
    \end{itemize}
    Let $X = [0, 1]$ and $S = [0, 1)$
    \begin{itemize}
        \item $x = 1$, $x = 0$, $x = 0.5$?
        \item What is $\interior S$? What is $\partial S$?
    \end{itemize}
\end{frame}

\begin{frame}[allowframebreaks]
    \frametitle{Interior, Exterior, Boundary, and Closure}
    \begin{definition}
        \begin{itemize}
            \item A point $x \in X$ is a \textbf{limit point} of $S$ if every
            open ball centered at $x$ contains at least one element of $S$
            other than $x$.
            \item A point $x \in S$ is an \textbf{isolated point} of $S$ if
            some open ball centered at $x$ contains no element of $S$ other
            than $x$.
            \item The \textbf{closure} of $S$ is the union of $S$ and the set
            of limit points of $S$, and is denoted by $\bar S$.
        \end{itemize}
    \end{definition}
    \framebreak
    \begin{proposition}
        \begin{enumerate}
            \item If $x$ is a limit point of $S$, then every open ball
            centered at $x$ contains an infinite number of points from $S$
            \item $\bar S = (\text{ext}\, S)^c$
            \item $\bar S = \interior S \cup \partial S$ 
        \end{enumerate}
    \end{proposition}
    Prove some using the definitions
\end{frame}

\begin{frame}
    \frametitle{Interior, Exterior, Boundary, and Closure}
    \framesubtitle{Examples}
    \begin{itemize}
        \item If $S = \{1/n: n \in \NN\} \subset \RR$, then every point in
        $S$ is an isolated point. The only limit point is 0.
        \item Consider $\QQ \subset \RR$. Then $\interior \QQ = \emptyset$,
        $\partial \QQ = \RR$, and $\bar \QQ = \RR$.
    \end{itemize}
\end{frame}

\begin{frame}[allowframebreaks]
    \frametitle{Open and Closed Sets}
    A set is open if all its elements are interior points.
    \begin{definition}
        \begin{itemize}
            \item A set $S \subset X$ is \textbf{open} iff $S \subset \interior S$.
            \item A set $S \subset X$ is \textbf{closed} iff $S^c$ is open.
        \end{itemize}
    \end{definition}
        \begin{proposition}
        \begin{enumerate}
            \item Any open ball is an open set
            \item A set $S$ is closed iff $S = \bar S$
            \item The sets $\bar S$ and $\partial S$ are closed
        \end{enumerate}
    \end{proposition}
    \framebreak
    \begin{theorem}
        \begin{enumerate}
            \item Any union of open sets is open; any intersection of closed
            sets is closed.
            \item Any finite intersection of open sets is open; any finite
            union of closed sets is closed.
        \end{enumerate}
    \end{theorem}
    Example: $(0, 1) = \bigcup_{n=2}^\infty[1/n, 1 - 1/n]$ (prove this)
\end{frame}

\section{Sequences and Convergence}

\begin{frame}
    \frametitle{Sequences}
    \begin{definition}
        \begin{itemize}
            \item A \textbf{sequence} in a set $X$ is a function from $\NN$
            to $X$. We write $x_1, x_2, \ldots$, or $(x_n)_{n=1}^\infty$, or
            simply $(x_n)$ for the sequence.
            \item Given a sequence $(x_n)$, a \textbf{subsequence} is a
            sequence $(x_{n_i})$ where $(n_i)$ is a strictly increasing
            sequence in $\NN$.
        \end{itemize}        
    \end{definition}
    Examples: $x_n = 1/n$, $x_n = 2n + 1$, etc.
\end{frame}

\begin{frame}
    \frametitle{Convergence of Sequences}
    \begin{definition}
        A sequence $(x_n)$ in a metric space $(X, d)$ \textbf{converges} to a
        \textbf{limit} $x\in X$, written as $x_n \to x$ or $\lim_{n\to\infty}
        x_n = x$, if for every $\epsilon > 0$, there exists $N \in \NN$ such
        that
        \begin{equation*}
            n \geq N \implies d(x_n, x) < \epsilon.
        \end{equation*}
    \end{definition}
\end{frame}

\begin{frame}
    \frametitle{Convergence of Sequences}
    \framesubtitle{Examples}
    \begin{itemize}
        \item $(x_n) = (1, 1, \ldots)$
        \item $x_n = 1/n$
        \item $(x_1, x_2, \ldots) = (1, 1/2, 1, 1/3, \ldots)$
    \end{itemize}
\end{frame}

\begin{frame}[allowframebreaks]
    \frametitle{Convergence of Sequences}
    \begin{proposition}
        A sequence in a metric space can have at most one limit.
    \end{proposition}
    \emph{Proof:} Suppose $x_n \to x$ and $x_n \to y$ as $n \to \infty$.
    Suppose $x \neq y$ and let $d(x, y) = r > 0$.

    There exists $N_1, N_2 \in \NN$ such that
    \begin{gather*}
        n \geq N_1 \implies d(x_n, x) < r/2\\
        n \geq N_2 \implies d(x_n, y) < r/2
    \end{gather*}
    Let $N = \max\{N_1, N_2\}$. Then $d(x, y) \leq d(x, x_N) + d(y, x_N) <
    r$, which is a contradiction.

    (Bounding the distance between two points using triangle inequality is a
    very important method that will be used over and over again.)

    \framebreak
    \begin{definition}
        A sequence $(x_n)$ is \textbf{bounded} if the set $\{x_n: n\in\NN\}$
        is bounded.
    \end{definition}
    \begin{proposition}
        A convergent sequence in a metric space is bounded.
    \end{proposition}
    \emph{Proof:} Let $x_n \to x$. There exists an $N \in \NN$ such that
    $d(x_n, x) < 1$ for all $n \geq N$. Let
    \begin{equation*}
        M = \max\{d(x_1, x), d(x_2, x), \ldots, d(x_N, x), 1\}.
    \end{equation*}
    Then $d(x_n, x) \leq M$ for all $n \in \NN$.

    (Using convergence to handle the ``tail'' of a sequence is an important
    method.)
\end{frame}

\begin{frame}
    \frametitle{Sequences and the Closure of a Set}
    \begin{theorem}
        Let $S$ be a subset of $X$. Then $x \in \bar S$ iff there is a
        sequence $(x_n) \subset S$ such that $x_n \to x$.
    \end{theorem}
    \begin{corollary}
        Let $S$ be a subset of $X$. Then $S$ is closed iff
        \begin{equation*}
            (x_n) \subset S \text{ and } x_n \to x \implies x \in S.
        \end{equation*}
    \end{corollary}
    \emph{Proof:} Suppose $(x_n) \subset S$ and $x_n \to x$. If $x_n=x$ for
    some $n$, then $x\in S\subset\bar S$. Otherwise, every open ball centered
    at $x$ contains a term $x_n\in S$ other than $x$, so $x$ is a limit point
    of $S$ and hence $x\in\bar S$.

    Conversely, suppose $x \in \bar S$. If $x\in S$, take the constant
    sequence $x_n=x$. If $x\notin S$, choose $x_n\in B_{1/n}(x)\cap S$ for
    each $n\in\NN$. Then $d(x_n,x)<1/n$, so $x_n\to x$.
\end{frame}


\begin{frame}[allowframebreaks]
    \frametitle{Sequences in $\RR$}
    \begin{proposition}
        Suppose $x_n \to x$ and $y_n \to y$. Then
        \begin{enumerate}
            \item $x_n + y_n \to x + y$
            \item $c x_n \to cx$ for all $c \in \RR$
            \item $x_n y_n \to xy$
            \item $1/x_n \to 1/x$ if $x_n \neq 0$ and $x \neq 0$
        \end{enumerate}
    \end{proposition}
    \begin{proposition}
        \begin{enumerate}
            \item If $0 \leq x_n \leq y_n$ for all $n \geq N$, and $y_n \to
            0$, then $x_n \to 0$.
            \item If $x_n \leq y_n$ for all $n \geq N$, $x_n \to x$, and $y_n
            \to y$, then $x \leq y$.
        \end{enumerate}
    \end{proposition}
    \framebreak
    \begin{theorem}
        Every bounded monotone sequence in $\RR$ has a limit in $\RR$.
    \end{theorem}
    The theorem can be proved using the Completeness Axiom.
\end{frame}

\section{Completeness}

\begin{frame}[allowframebreaks]
    \frametitle{Cauchy Sequences}
    \begin{itemize}
        \item The definition for convergence requires that we know the limit
        a priori.
        \item We would like a criterion for convergence that does not depend
        on the limit itself.
    \end{itemize}
    \begin{definition}
        Let $(x_n) \subset X$. Then $(x_n)$ is a \textbf{Cauchy sequence} if
        for every $\epsilon > 0$, there exists $N \in \NN$ such that
        \begin{equation*}
            m, n \geq N \implies d(x_m, x_n) < \epsilon.
        \end{equation*}
    \end{definition}
    \framebreak
    \begin{proposition}
        \begin{enumerate}
            \item If a sequence is Cauchy, then it is bounded.
            \item If a sequence converges, then it is Cauchy.
        \end{enumerate}
    \end{proposition}
    Similar to convergent sequences, Cauchy sequences are also bounded, but
    the second result above suggests that Cauchy is weaker.

    When is a Cauchy sequence convergent?
\end{frame}

\begin{frame}[allowframebreaks]
    \frametitle{Completeness}
    \begin{definition}
        A metric space $(X, d)$ is \textbf{complete} if every Cauchy sequence
        in $X$ has a limit in $X$.
    \end{definition}

    \begin{proposition}
        $\RR^n$ with the Euclidean metric is a complete metric space.
    \end{proposition}
    \framebreak
    \begin{proposition}
        Let $(X, d)$ be a metric space and let $S \subset X$.
        \begin{enumerate}
            \item If $(S, d)$ is complete, then $S$ is closed.
            \item If $(X, d)$ is complete, then $S$ is closed in $X$ iff $(S,
            d)$ is complete.
        \end{enumerate}
    \end{proposition}
    \emph{Proof:} 1. To prove that $S$ is closed, we pick any sequence $x_n
    \subset S$ converging to $x \in X$. Then $(x_n)$ is Cauchy in $S$ and
    thus has a limit in $S$.

    2. Suppose $S$ is closed. Since any Cauchy sequence in $S$ is a Cauchy
    sequence in $X$, and thus has a limit in $X$. Since $S$ is closed, any
    sequence $(x_n) \subset S$ that converges has a limit in $S$. Therefore,
    $(S, d)$ is complete. The other direction follows from 1.
\end{frame}

\begin{frame}[allowframebreaks]
    \frametitle{Contraction Mapping Theorem}
    \begin{definition}
        Let $(X, d)$ be a metric space. We say $F: S\subset X \to X$ is a
        \textbf{contraction} if there exists $0 \leq \lambda < 1$ such that
        \begin{equation*}
            d(F(x), F(y)) \leq \lambda d(x, y)
        \end{equation*}
        for all $x, y \in S$.
    \end{definition}
    \begin{definition}
        We say $x^*\in S$ is a \textbf{fixed} point of $F: S\subset X \to X$
        if $F(x^*) = x^*$.
    \end{definition}
    \framebreak
    \begin{theorem}[Contraction Mapping Theorem]
        Let $(X, d)$ be a \emph{complete} metric space and let $F: X \to X$
        be a contraction. Then $F$ has a unique fixed point $x^*$ and $F^n(x)
        \to x^*$ as $n \to \infty$ for all $x \in X$.
    \end{theorem}
    \emph{Notes:} here $F^n(x)$ means apply $F$ iteratively $n$ times:
    \begin{equation*}
        x_1 = F(x),\, x_2 = F(x_1),\, x_3 = F(x_2),\, \ldots,\, x_n = F(x_{n-1})
    \end{equation*}
    \framebreak
    \begin{itemize}
        \item In economics, the contraction mapping theorem is mainly used in
        dynamic programming
        \item The Bellman operator is a contraction mapping 
    \end{itemize}
\end{frame}

\section{Compactness}

\begin{frame}[allowframebreaks]
    \frametitle{Subsequences}
    \begin{theorem}
        If a sequence in a metric space converges, then every subsequence
        converges to the same limit as the original sequence.
    \end{theorem}
    \begin{theorem}
        If a Cauchy sequence $(x_n)$ in a metric space has a subsequence
        converging to $x$, then $x_n \to x$.
    \end{theorem}
    \emph{Proof:} The first is by applying the definition. To prove the
    second theorem, pick an arbitrary $\epsilon > 0$. There exists $N \in
    \NN$ such that $d(x_m, x_n) < \epsilon/2$ for all $m, n \geq N$. If a
    subsequence $(x_{n_i})$ converges to $x$, then there exists $M \in \NN$
    such that $d(x_{n_i}, x) < \epsilon/2$ for all $i \geq M$. Choose one
    $j\geq M$ such that $n_j\geq N$. Then for every $m\geq N$,
    \begin{equation*}
        d(x, x_m) \leq d(x, x_{n_j}) + d(x_{n_j}, x_m) < \epsilon.
    \end{equation*}
    \framebreak
    \begin{theorem}[Bolzano-Weierstrass Theorem]
        Every bounded sequence in $\RR^n$ has a convergent subsequence.
    \end{theorem}
    The same result is not true for an arbitrary metric space!

    \begin{corollary}
        If $S \subset \RR^n$, then $S$ is closed and bounded iff every
        sequence from $S$ has a subsequence which converges to a limit in
        $S$.
    \end{corollary}
    \emph{Proof:} First suppose $S$ is closed and bounded. Since $S$ is
    bounded, every sequence has a convergent subsequence. Since $S$ is
    closed, the limit is in $S$.

    For the other direction, closedness is obvious. We prove $S$ is bounded
    by contradiction. Suppose $S$ is unbounded. Then for every $n \in \NN$,
    there exists $x_n \in S$ such that $|x_n| > n$. Any subsequence of
    $(x_n)$ is unbounded and thus cannot converge.
\end{frame}

\begin{frame}[allowframebreaks]
    \frametitle{Compactness}
    \begin{definition}
        A subset $S$ of a metric space $(X, d)$ is \textbf{compact} if every
        sequence from $S$ has a subsequence that converges to a limit in $S$.
    \end{definition}
    This definition is actually called \emph{sequential compactness}, but it
    is equivalent to compactness for metric spaces.

    The Bolzano-Weierstrass Theorem implies the following theorem (although
    it was proved independently).
    \begin{theorem}[Heine-Borel Theorem]
        A subset of $\RR^n$ is compact iff it is closed and bounded.
    \end{theorem}
    \framebreak
    \begin{proposition}
        If a subset of a metric space $(X, d)$ is compact, then it is bounded
        and closed (in $X$).
    \end{proposition}
    \emph{Proof:} Suppose $S \subset X$ is compact. Then for any sequence
    $(x_n) \subset S$ that converges to $x \in X$, there exists a subsequence
    that converges to a limit in $S$. By a previous theorem, the limit must
    be $x$. Hence $S$ is closed.

    Suppose $S$ is not bounded. Then we can build a sequence that has no
    convergent subsequence.
    \framebreak
    \begin{proposition}
        Let $S$ be a subset of a compact metric space $(X, d)$. Then $S$ is
        compact iff $S$ is closed in $X$.
    \end{proposition}
    \emph{Proof:} If $S$ is compact, then it is closed. Suppose that $S$ is
    closed. Since $X$ is compact, any sequence in $S \subset X$ has a
    convergent subsequence in $X$. Since $S$ is closed, the limit is also in
    $S$. Hence $S$ is compact.
\end{frame}

\section{Continuity}

\begin{frame}[allowframebreaks]
    \frametitle{Limits of Functions}
    \begin{definition}
        Let $f:(A \subset X) \to Y$ where $X$ and $Y$ are metric spaces and
        let $a$ be a limit point of $A$. If for every sequence $(x_n) \subset
        A\setminus\{a\}$, $x_n \to a$ implies $f(x_n) \to b$, then we say
        $f(x)$ approaches $b$ as $x$ approaches $a$, or $f$ has
        \textbf{limit} $b$ at $a$, and write
        \begin{equation*}
            f(x) \to b \text{ as } x \to a,
        \end{equation*}
        or
        \begin{equation*}
            \lim_{x\to a}f(x) = b.
        \end{equation*}
    \end{definition}
    \framebreak
    \begin{definition}
        If $X = \RR$ and $A$ is an interval with $a$ as left or right
        endpoint, then we write
        \begin{equation*}
            \lim_{x\to a^+} f(x) \text{ or } \lim_{x\to a^-} f(x)
        \end{equation*}
        and say the limit as $x$ approaches from the right or the left,
        respectively.
    \end{definition}
    Examples:
    \begin{itemize}
        \item Let $A = (-\infty, 0) \cup (0, \infty)$. Let $f(x) = 1$ for $x
        \in A$. Then $\lim_{x \to 0}f(x) = 1$.
        \item Draw some diagrams
    \end{itemize}
    \framebreak
    We have an equivalent definition of a limit using the usual
    $\epsilon$-$\delta$ language
    \begin{theorem}
        Suppose $(X, d)$ and $(Y, \rho)$ are metric spaces, $A \subset X$,
        $f: A \to Y$, and $a$ is a limit point of $A$. Then the following are
        equivalent
        \begin{enumerate}
            \item $\lim_{x \to a} f(x) = b$
            \item For every $\epsilon > 0$, there exists a $\delta > 0$ such
            that
            \begin{equation*}
                x \in A \setminus \{a\} \text{ and } d(x, a) < \delta
                \implies \rho(f(x), b) < \epsilon
            \end{equation*}
        \end{enumerate}
    \end{theorem}
\end{frame}

\begin{frame}[allowframebreaks]
    \frametitle{Continuity}
    \begin{definition}
        Let $f: A(\subset X) \to Y$ where $X$ and $Y$ are metric spaces and
        let $a \in A$. We say $f$ is \textbf{continuous} at $a$ if $a$ is an
        isolated point of $A$ or if $a$ is a limit point and $\lim_{x \to a}
        f(x) = f(a)$.

        If $f$ is continuous at every $a \in A$, then we say $f$ is
        \textbf{continuous}. The set of all such continuous functions is
        denoted by $\cC(A; Y)$ or $\cC(A)$ if $Y = \RR$.
    \end{definition}
    \framebreak
    We also have equivalent definitions of continuity
    \begin{proposition}
        Suppose $(X, d)$ and $(Y, \rho)$ are metric spaces, $A \subset X$,
        $f: A \to Y$, and $a \in A$. Then the following are equivalent:
        \begin{enumerate}
            \item $f$ is continuous at $a$
            \item whenever $(x_n) \subset A$ and $x_n \to a$, then $f(x_n)
            \to f(a)$
            \item for each $\epsilon > 0$, there exists $\delta > 0$ such
            that
            \begin{equation*}
                x\in A \text{ and } d(x, a) < \delta \implies \rho\left(f(x),
                f(a) \right) < \epsilon
            \end{equation*}
        \end{enumerate}
    \end{proposition}
\end{frame}

\begin{frame}
    \frametitle{Continuity}
    \framesubtitle{Examples}
    \begin{itemize}
        \item Show that $f(x) = x$ is continuous at $x = 1$
        \item Show that $f(x) = x^2$ is continuous at every point in $\RR$
    \end{itemize}
\end{frame}

\begin{frame}[allowframebreaks]
    \frametitle{Continuity}
    \begin{theorem}
        \begin{enumerate}
            \item Let $f$ and $g$ be two real-valued functions. If $f$ and
            $g$ are continuous at $a$, then $f + g$, $fg$, and $f/g$ (if
            $g(a) \neq 0$) are continuous at $a$
            \item If $f$ is continuous at $a$ and $g$ is continuous at
            $f(a)$, then $g \circ f$ is continuous at $a$
        \end{enumerate}
    \end{theorem}
    \begin{theorem}[Intermediate Value Theorem]
        Let $f$ be a continuous real function on the interval $[a, b]$. If
        $f(a) < f(b)$ and $c$ is such that $f(a) < c < f(b)$, then there
        exists a point $x \in (a, b)$ such that $f(x) = c$.
    \end{theorem}
    \framebreak
    \begin{theorem}
        Let $f: X \to Y$ where $(X, d)$ and $(Y, \rho)$ are metric spaces.
        Then the following are equivalent
        \begin{enumerate}
            \item $f$ is continuous
            \item $f^{-1}(E)$ is open in $X$ whenever $E$ is open in $Y$
            \item $f^{-1}(C)$ is closed in $X$ whenever $C$ is closed in $Y$
        \end{enumerate}
    \end{theorem}
    The continuous image of an open set is not necessarily open and the
    continuous image of a closed set is not necessarily closed.

    For example, let $f(x) = x^2$ and then $f((-1, 1)) = [0, 1)$; let $f(x) =
    e^x$ and $f(\RR) = (0, \infty)$
\end{frame}

\begin{frame}
    \frametitle{Continuous Functions on Compact Sets}
    \begin{theorem}
        Let $f:K(\subset X) \to Y$ be a continuous function, where $X$ and
        $Y$ are metric spaces, and $K$ is compact. Then
        \begin{enumerate}
            \item $f(K)$ is compact
            \item If $Y \subset \RR$ and $K$ is nonempty, then $f$ is bounded
            (above and below) and has a maximum and a minimum
        \end{enumerate}
    \end{theorem}
    Examples:
    \begin{itemize}
        \item Let $f(x) = 1/x$ for $x \in (0, 1]$
        \item Let $f(x) = x$ for $x \in [0, 1)$
        \item Let $f(x) = 1/x$ for $x \in [1, \infty)$
    \end{itemize}
\end{frame}

\begin{frame}
    \frametitle{The Space of Bounded Continuous Functions}
    \begin{definition}
        The set of \emph{bounded} continuous functions $f: A \to Y$ is
        denoted by $\bB\cC(A; Y)$.
    \end{definition}
    \begin{theorem}
        Suppose $A \subset X$, $(X, d)$ is a metric space and $(Y, \rho)$ is
        a complete metric space. Then $\bB\cC(A; Y)$ is a complete metric
        space under the \textbf{uniform metric} $d_u$ defined by
        \begin{equation*}
            d_u(f, g) := \sup_{x \in A}\rho\Big(f(x), g(x)\Big).
        \end{equation*}
    \end{theorem}
    Several important results:
    \begin{enumerate}
        \item $(\bB\cC(A; Y), d_u)$ is a metric space
        \item $\bB\cC(A; Y)$ is closed under uniform convergence: if $f_n \to
        f$ in the uniform metric, then $f$ is also bounded continuous
    \end{enumerate}
\end{frame}

\end{document}